% ============================================================
%  TEMPLATE DE TESIS - Universidad Nacional del Sur
%  Licenciatura en Ciencias de la Computación
%  Alumna: Candela Ledesma
% ============================================================

\documentclass[12pt, a4paper]{report}

% ---- Paquetes ----
\usepackage[utf8]{inputenc}
\usepackage[T1]{fontenc}
\usepackage[spanish]{babel}
\usepackage{geometry}
\usepackage{graphicx}
\usepackage{hyperref}
\usepackage{setspace}
\usepackage{fancyhdr}
\usepackage{titlesec}
\usepackage{tocloft}
\usepackage{caption}
\usepackage{subcaption}
\usepackage{listings}
\usepackage{xcolor}
\usepackage{booktabs}
\usepackage{amsmath}
\usepackage{amssymb}
\usepackage{csquotes}
\usepackage{enumitem}
\usepackage{parskip}
\usepackage{emptypage}
\usepackage{microtype}
\usepackage{float}      % permite [H] para fijar figuras exactamente en su posición

% ---- Márgenes ----
\geometry{
  top=2.5cm,
  bottom=2.5cm,
  left=3cm,
  right=2.5cm
}

% ---- Interlineado ----
\onehalfspacing

% ---- Colores para código ----
\definecolor{codegreen}{rgb}{0,0.6,0}
\definecolor{codegray}{rgb}{0.5,0.5,0.5}
\definecolor{codepurple}{rgb}{0.58,0,0.82}
\definecolor{backcolour}{rgb}{0.97,0.97,0.97}

% ---- Estilo para bloques de código ----
\lstdefinestyle{mystyle}{
  backgroundcolor=\color{backcolour},
  commentstyle=\color{codegreen},
  keywordstyle=\color{codepurple}\bfseries,
  numberstyle=\tiny\color{codegray},
  stringstyle=\color{codegreen},
  basicstyle=\ttfamily\footnotesize,
  breakatwhitespace=false,
  breaklines=true,
  captionpos=b,
  keepspaces=true,
  numbers=left,
  numbersep=5pt,
  showspaces=false,
  showstringspaces=false,
  showtabs=false,
  tabsize=2,
  frame=single,
  rulecolor=\color{codegray}
}
\lstset{style=mystyle}

% ---- Estilo de encabezado y pie de página ----
\pagestyle{fancy}
\fancyhf{}
\fancyhead[L]{\small\nouppercase{\leftmark}}
\fancyhead[R]{\small Candela Ledesma}
\fancyfoot[C]{\thepage}
\renewcommand{\headrulewidth}{0.4pt}
\renewcommand{\chaptermark}[1]{\markboth{\thechapter.\ #1}{}}

% ---- Estilo de capítulos ----
\titleformat{\chapter}[display]
  {\normalfont\Large\bfseries}
  {\chaptertitlename\ \thechapter}{10pt}{\huge}
\titlespacing*{\chapter}{0pt}{-10pt}{20pt}

% ---- Hipervínculos (siempre al final del preámbulo) ----
\hypersetup{
  colorlinks=true,
  linkcolor=black,
  urlcolor=blue,
  citecolor=black,
  pdftitle={Título de la Tesis - Candela Ledesma},
  pdfauthor={Candela Ledesma}
}

% ============================================================
\begin{document}

% ============================================================
%  PORTADA
% ============================================================
\begin{titlepage}
  \centering

  % \includegraphics[width=0.25\textwidth]{logo_uns.png}\\[1cm]

  \vspace{1cm}

  {\LARGE \textbf{Documento de Tesis}}\\[0.3cm]
  {\large Licenciatura en Ciencias de la Computación}\\[2cm]

  {\huge \textbf{[Título de la Tesis Placeholder]}}\\[0.5cm]
  {\large Subtítulo descriptivo del proyecto (opcional)}\\[2.5cm]

  \begin{flushleft}
    \large
    \textbf{Alumna:} Candela Ledesma \\[0.4cm]
    \textbf{Director:} Dr.\ Sebastian Gottifredi \\[0.4cm]
    \textbf{Departamento:} Ciencias e Ingeniería de la Computación \\[0.4cm]
    \textbf{Universidad:} Universidad Nacional del Sur \\[0.4cm]
    \textbf{Fecha:} \the\year
  \end{flushleft}

  \vfill
\end{titlepage}

% ============================================================
%  PÁGINAS PRELIMINARES
% ============================================================
\pagenumbering{roman}
\setcounter{page}{1}

\newpage
\thispagestyle{empty}
\mbox{}

% ---- Índice general ----
\newpage
\tableofcontents

% ============================================================
%  CUERPO PRINCIPAL
% ============================================================
\newpage
\pagenumbering{arabic}
\setcounter{page}{1}

% ============================================================
%  CAPÍTULO 1 - INTRODUCCIÓN
% ============================================================
\chapter{Introducción}
\label{cap:introduccion}

\section{Motivación}

Las universidades publican sus planes de estudio en formato PDF: documentos diseñados para
ser leídos por personas, no procesados por sistemas. Para un estudiante que quiere entender
qué materias debe cursar antes de inscribirse a otra, o planificar su recorrido académico a
lo largo de años, ese PDF ofrece poca ayuda. La información está ahí, pero no es
interactiva, no responde preguntas, y no permite explorar dependencias entre materias de
forma visual.

Este problema existe en la Universidad Nacional del Sur (UNS), donde cada carrera tiene su
propio plan de estudios con correlativas, agrupadores de optativas, requisitos especiales
expresados en prosa y estructuras que varían entre departamentos. Un estudiante que quiere
saber si puede inscribirse a una materia debe rastrear manualmente una cadena de
dependencias que puede atravesar varios años del plan.

La motivación de este trabajo es resolver ese problema: transformar los planes de estudio
estáticos de la UNS en una aplicación web interactiva que permita a cada estudiante
visualizar su situación académica, marcar el progreso de sus materias y entender qué
caminos tiene disponibles para avanzar en su carrera.

\section{Problema y Desafíos}

El punto de partida del problema es la brecha entre el formato de publicación y el formato
que requiere una aplicación interactiva. Los PDFs de planes de estudio de la UNS no siguen
una estructura uniforme: algunas carreras usan tablas, otras listas, algunas expresan
correlativas como IDs numéricos, otras las describen en prosa. Esta heterogeneidad hace que
cualquier intento de extracción automática enfrente casos borde difíciles de anticipar.

El primer desafío técnico es convertir esos documentos en un JSON estructurado y validado
que capture fielmente la información del plan: materias, años, cuatrimestres, correlativas
para cursar y para rendir, agrupadores de optativas, grupos de idioma y requisitos
especiales. Cualquier error en esta conversión se propaga directamente a la experiencia del
usuario: si una correlativa está mal registrada, el sistema puede habilitar materias que el
estudiante no debería poder cursar, o bloquear materias que sí debería tener disponibles.

El segundo desafío es la escala del problema de conversión. La UNS ofrece decenas de
carreras. Mantener los JSONs actualizados manualmente cada vez que un plan cambia no es
viable. Se necesita un pipeline que reduzca la intervención humana al mínimo y que pueda
incorporar nuevos planes con el menor esfuerzo posible.

El tercer desafío es la validación. Dado que el pipeline involucra modelos de inteligencia
artificial generativa, los outputs no son deterministas ni siempre correctos. Se requiere
una estrategia de evaluación que permita medir la calidad del JSON generado y detectar
errores antes de que lleguen a los usuarios.

\section{Objetivo General}

Diseñar e implementar un sistema que convierta planes de estudio universitarios en PDF a
datos estructurados, y los publique en una aplicación web que permita a los estudiantes
gestionar su progreso académico y explorar las dependencias entre materias de su carrera.

\section{Objetivos Específicos}

\begin{enumerate}
  \item Diseñar un schema JSON validado que represente fielmente la estructura de un plan de
        estudios, incluyendo correlativas, optativas e idiomas.
  \item Desarrollar un parser local en Python capaz de extraer materias, correlativas y
        agrupadores desde PDFs de planes de estudio de la UNS.
  \item Integrar un modelo de lenguaje de gran escala (Gemini) para generar JSONs
        directamente desde PDFs mediante visión nativa, como alternativa al parser local.
  \item Implementar una métrica de similitud que permita comparar el output del modelo
        contra el ground truth del parser local.
  \item Desarrollar una aplicación web que permita a los estudiantes visualizar planes de
        estudio, marcar su progreso y explorar correlativas de forma interactiva.
  \item Implementar un panel de administración que permita subir PDFs, comparar outputs del
        parser y del modelo, y publicar planes sin intervención técnica.
\end{enumerate}

\section{Organización del documento}

El capítulo~\ref{cap:contexto} describe el contexto institucional y las tecnologías
utilizadas. El capítulo~\ref{cap:arquitectura} presenta la arquitectura de la solución y el
modelo de datos. El capítulo~\ref{cap:implementacion} detalla la implementación de los
módulos principales. El capítulo~\ref{cap:evaluacion} evalúa los resultados obtenidos
mediante la métrica de comparación entre el parser local y el modelo de IA. El
capítulo~\ref{cap:conclusiones} presenta las conclusiones, lecciones aprendidas y trabajo
futuro.

% ============================================================
%  CAPÍTULO 2 - CONTEXTO Y TECNOLOGÍAS
% ============================================================
\chapter{Contexto y Tecnologías}
\label{cap:contexto}

\section{Contexto}

La Universidad Nacional del Sur (UNS) es una universidad pública argentina con sede en
Bahía Blanca. Ofrece carreras de grado distribuidas en múltiples departamentos, cada uno
con autonomía para definir la estructura de sus planes de estudio. Esta autonomía se refleja
en la heterogeneidad de los documentos: algunas carreras organizan el plan en tablas por
año y cuatrimestre, otras en listas lineales; algunas expresan las correlativas como
identificadores numéricos, otras las describen en prosa. No existe un formato unificado.

Los planes de estudio se publican únicamente en formato PDF, un formato diseñado para
impresión y lectura humana. Para un sistema informático, extraer información estructurada de
esos documentos requiere resolver primero el problema de la conversión: transformar un
documento orientado a la presentación en datos procesables por una aplicación.

Este trabajo se enmarca también en el surgimiento de los modelos de lenguaje de gran escala
(LLMs) con capacidades de visión: modelos que pueden recibir un PDF como imagen y razonar
sobre su contenido visual sin requerir extracción de texto previa. Esta capacidad abre una
vía alternativa al parsing clásico basado en reglas, con la posibilidad de generalizar a
documentos con estructuras arbitrarias que serían difíciles de cubrir con heurísticas
escritas a mano.

El sistema desarrollado en este trabajo aprovecha ambos enfoques: un parser local en Python
como fuente de verdad verificable, y un modelo de IA generativa (Gemini) como mecanismo de
conversión escalable, con un panel de administración que permite comparar ambos outputs
antes de publicar.

\section{Herramientas y Tecnologías}

A continuación se describen las principales tecnologías utilizadas en el desarrollo del
proyecto.

\subsection{Prisma y PostgreSQL}

La base de datos relacional es PostgreSQL, accedida a través de Prisma ORM v6. Prisma
cumple tres roles en el proyecto: define el schema de la base de datos en un lenguaje
declarativo propio (\texttt{schema.prisma}), genera un cliente TypeScript con tipado
completo a partir de ese schema, y gestiona las migraciones de base de datos.

La instancia de PostgreSQL se aloja en Neon, una plataforma de PostgreSQL serverless con
soporte a conexiones HTTP y WebSocket, elegida por su integración nativa con el modelo de
despliegue serverless de Vercel. Neon provee dos URLs de conexión: una pooled (vía
PgBouncer) para las funciones serverless, y una directa para las migraciones de Prisma que
requieren una conexión persistente.

\subsection{Gemini y la API de Google GenAI}

El componente de inteligencia artificial usa la API de Google Generative AI
(\texttt{@google/genai}) para enviar PDFs directamente al modelo Gemini~2.5~Flash y recibir
un JSON estructurado como respuesta. El PDF se transmite como \texttt{inlineData} en base64
con MIME type \texttt{application/pdf}, lo que permite que Gemini procese el documento con
sus capacidades de visión nativa sin extracción de texto previa.

El modelo opera con temperatura~0 para maximizar el determinismo de los outputs. El prompt
del sistema (almacenado en \texttt{frontend/lib/ai/prompt.ts} y editable desde la tabla
\texttt{AdminConfig} de la base de datos sin necesidad de redeploy) instruye al modelo sobre
el schema JSON de salida, las reglas de agrupadores de optativas, los grupos de idioma y los
casos borde más frecuentes. El sistema soporta cuatro modelos intercambiables desde el panel
admin: \texttt{gemini-2.5-flash} (modelo principal), \texttt{gemini-2.5-flash-lite},
\texttt{gemini-2.5-pro} y \texttt{gemma-4-26b-a4b-it}.

\subsection{Parser local en Python}

El parser local está implementado en Python y expuesto como una API REST con FastAPI y
Uvicorn. Recibe un PDF por multipart form-data y devuelve el JSON del plan de estudios.
Internamente usa \texttt{pdfplumber} para la extracción de texto y la detección de layout,
y aplica reglas de parsing específicas para los formatos de la UNS.

Este parser no reemplaza al modelo de IA: cumple el rol de generador del ground truth,
produciendo el JSON de referencia contra el cual se evalúa la calidad del output de Gemini.
La comparación entre ambos outputs se realiza con un script de evaluación
(\texttt{scripts/comparar\_json.py}) que calcula métricas de similitud sobre materias,
correlativas y agrupadores.

\subsection{Despliegue: Vercel y Render}

El frontend Next.js se despliega en Vercel. Vercel compila y publica automáticamente cada
push a la rama principal, genera previews por pull request, y sirve las funciones
serverless con routing optimizado para Next.js. Los deployments de preview usan branches
de Neon aisladas, sin contaminar la base de datos de producción.

El parser local en Python se despliega en Render como servicio independiente. Render
permite workers persistentes de larga duración, necesarios para Uvicorn: el parser no
puede correr como función serverless porque el modelo stateless de Vercel es incompatible
con el tiempo de startup de un servidor Python con dependencias como \texttt{pdfplumber}.
Esta separación implica que el panel admin realiza dos requests en paralelo al subir un
PDF: uno a Vercel (para Gemini) y otro a Render (para el parser local).

\subsection{Testing}

La suite de tests cubre tres niveles. Los \textbf{tests unitarios del parser} en Python
(pytest) verifican que las reglas de parsing produzcan el JSON correcto para cada PDF del
dataset. Los \textbf{tests de integración del frontend} (Vitest) cubren la lógica de
validación de schemas, sincronización de progreso y utilidades de correlativas. Los
\textbf{tests end-to-end} (Playwright) verifican los flujos principales de la aplicación
web, incluyendo navegación, marcado de progreso y acceso al panel de administración.

\section{Comparación de alternativas}
\label{sec:comparacion-alternativas}

La tabla~\ref{tab:alternativas} resume las principales decisiones tecnológicas del sistema
y el motivo de cada elección.

\begin{table}[htbp]
\centering
\small
\begin{tabular}{p{3.2cm} p{3.2cm} p{7cm}}
\toprule
\textbf{Elegido} & \textbf{Descartado} & \textbf{Motivo principal} \\
\midrule
Parser + Gemini (híbrido)
  & Solo parser o solo IA
  & El parser garantiza ground truth verificable; Gemini escala a formatos arbitrarios.
    Ninguno es suficiente solo. \\
\addlinespace
PostgreSQL (Prisma)
  & MongoDB u otra BD documental
  & Las entidades del sistema son relacionales; solo el progreso del usuario tiene naturaleza
    documental y se serializa como JSON en un campo texto
    (véase sección~\ref{sec:modulo-progreso}). \\
\addlinespace
NextAuth v4
  & OAuth propio
  & Encapsula el flujo completo de Google OAuth en pocas líneas y se integra nativamente
    con Prisma. Suficiente para el alcance del sistema. \\
\addlinespace
Vercel + Neon
  & Servidor dedicado
  & Despliegue continuo sin configuración de infraestructura; escalado automático adecuado
    para la carga variable de un sistema universitario. \\
\bottomrule
\end{tabular}
\caption{Resumen de decisiones tecnológicas principales.}
\label{tab:alternativas}
\end{table}

La decisión más relevante es la arquitectura híbrida parser/IA. Un parser basado en reglas
ofrece alta predecibilidad y resultados verificables, pero requiere mantenimiento
proporcional a la diversidad de formatos de cada carrera. Un modelo de IA generativa
generaliza a formatos arbitrarios sin código adicional, pero sus outputs no son
deterministas. El sistema combina ambos: el parser genera el ground truth y Gemini opera
como motor de conversión escalable, con el panel admin como punto de validación entre
ambos.

% ============================================================
%  CAPÍTULO 3 - ARQUITECTURA DE LA SOLUCIÓN
% ============================================================
\chapter{Arquitectura de la Solución}
\label{cap:arquitectura}

\section{Descripción general}

La solución propuesta se apoya en una arquitectura [tipo de arquitectura, por ejemplo:
modular, cliente-servidor, etc.] que permite [ventaja principal].

\section{Diagrama de arquitectura}

% \begin{figure}[htbp]
%   \centering
%   \includegraphics[width=0.8\textwidth]{diagrama_arquitectura.png}
%   \caption{Diagrama general de la arquitectura del sistema.}
%   \label{fig:arquitectura}
% \end{figure}

La arquitectura se puede dividir en las siguientes capas:

\begin{itemize}
  \item \textbf{Capa de presentación:} [descripción]
  \item \textbf{Capa de lógica de negocio:} [descripción]
  \item \textbf{Capa de datos:} [descripción]
\end{itemize}

\section{Modelo de datos}

El sistema utiliza PostgreSQL como base de datos relacional, accedida a través de Prisma
ORM. El schema está diseñado alrededor de tres ejes: autenticación y perfil de usuario,
gestión del ciclo de vida de los planes de estudio, y trazabilidad de acciones
administrativas.

\subsection{Entidades principales}

\subsubsection*{User}

Representa a un usuario autenticado. Se crea automáticamente al primer login con Google
(vía NextAuth). El campo \texttt{role} controla el nivel de acceso: \texttt{USER} solo
accede a la aplicación, \texttt{MODERATOR} puede enviar planes para revisión,
\texttt{ADMIN} puede publicar planes, gestionar usuarios y editar la configuración del
sistema.

\subsubsection*{Carrera y Departamento}

\texttt{Carrera} es la entidad central del dominio académico: tiene un identificador
\texttt{slug} único (ej: \texttt{ingenieria\_civil}), pertenece a un \texttt{Departamento}
(unidad académica, también identificada por slug), y puede tener múltiples versiones de
plan (\texttt{PlanVersion}).

\subsubsection*{PlanVersion}

Modela las versiones del plan de estudios de una carrera (ej: \texttt{v1}, \texttt{v2}).
Referencia el archivo JSON con el contenido del plan (\texttt{jsonFile}) y apunta al
\texttt{Plan} publicado que respalda esa versión.

\subsubsection*{Plan}

Tabla unificada para el ciclo de vida de un plan. Anteriormente existían tres tablas
separadas (\texttt{PlanBorrador}, \texttt{PlanPendiente}, \texttt{PlanPublicado}) con
campos inconsistentes entre sí. Se consolidaron en una sola tabla con un enum
\texttt{estado}:

\begin{table}[htbp]
\centering
\begin{tabular}{ll}
\toprule
\textbf{Estado} & \textbf{Significado} \\
\midrule
\texttt{BORRADOR}  & Generado al parsear un PDF, pendiente de revisión \\
\texttt{PENDIENTE} & Enviado por moderador, esperando aprobación de admin \\
\texttt{PUBLICADO} & Visible para los usuarios \\
\texttt{INACTIVO}  & Desactivado temporalmente sin perder historial \\
\bottomrule
\end{tabular}
\caption{Estados del ciclo de vida de un \texttt{Plan}.}
\label{tab:estados-plan}
\end{table}

El campo \texttt{fuente} (enum \texttt{PlanFuente}) indica si el JSON fue generado por el
parser local Python (\texttt{PARSER}), por Gemini con visión nativa (\texttt{GEMINI}), o
fusionado manualmente en el panel admin (\texttt{MERGED}). El campo \texttt{esBackup}
reemplaza el mecanismo anterior donde los backups se guardaban con el slug
\texttt{"\{slug\}\_v1\_backup"} dentro de la misma tabla.

La restricción de unicidad \texttt{@@unique([slug, fuente, estado])} es central para el
flujo de comparación del panel admin: permite que coexistan simultáneamente un plan de
fuente \texttt{PARSER} y uno de fuente \texttt{GEMINI} en estado \texttt{BORRADOR} para
la misma carrera, habilitando la comparación side-by-side entre ambos outputs antes de
decidir cuál publicar. Al mismo tiempo, la restricción garantiza que no pueda existir más
de un plan \texttt{PUBLICADO} por carrera y fuente, evitando inconsistencias en la versión
visible para los usuarios.

\subsubsection*{UserPlanProgress}

Almacena el progreso del usuario en un plan: qué materias aprobó, cursó o no cursó. El
estado se representa como un objeto JSON serializado
\texttt{\{[materiaId]: "aprobada"\,|\,"cursada"\,|\,"no\_cursada"\}} que se sobreescribe
completo en cada sincronización. La justificación técnica de esta decisión de diseño se
desarrolla en la sección~\ref{sec:modulo-progreso}.

\subsubsection*{CarreraSeleccionada}

Gestiona la relación muchos-a-muchos entre \texttt{User} y \texttt{Carrera}: cada fila
indica que un usuario eligió seguir una carrera dentro de la aplicación. No representa
una inscripción académica formal. La restricción \texttt{UNIQUE([userId, careerId])}
evita duplicados en el panel del usuario.

\subsubsection*{UserPreference}

Singleton por usuario (relación 1-a-1). Almacena \texttt{activeCareerId} (carrera
seleccionada actualmente en el dashboard) y los timestamps de onboarding
(\texttt{onboardingCompletedAt}, \texttt{onboardingDismissedAt}), que controlan si se
muestra el wizard de bienvenida en el próximo acceso.

\subsubsection*{UserRecentPlan}

Registra el último plan que cada usuario abrió por carrera, con un upsert en cada
navegación. Esta tabla habilita la funcionalidad ``continuar donde lo dejaste'': al
cargar la pantalla principal, el sistema consulta \texttt{UserRecentPlan} para mostrar
directamente la versión del plan visitada por última vez.

\subsubsection*{ProgressShare}

Snapshot inmutable del progreso de un usuario en un momento dado, identificado por un
token público. Se crea al compartir el avance y nunca se actualiza: el enlace siempre
muestra el estado exacto del instante en que se compartió.

\subsubsection*{ScheduleBlock}

Bloques del planificador semanal. \texttt{dia} va de \texttt{0} (lunes) a \texttt{6}
(domingo); los horarios se almacenan como enteros \texttt{HHMM} (ej: \texttt{830},
\texttt{1000}). El campo \texttt{materiaId} es opcional, permitiendo bloques libres sin
materia asociada. El campo \texttt{color} permite asignar un color hexadecimal a cada
bloque para distinguirlos visualmente en la grilla.

\subsubsection*{AdminConfig}

Singleton (\texttt{id = "singleton"}) que almacena los prompts enviados a Gemini para
parsear PDFs. Editable desde el panel admin sin necesidad de redeploy: si existe la fila,
el sistema la usa; si no, cae al valor compilado en \texttt{frontend/lib/ai/prompt.ts}.

\subsubsection*{AuditLog}

Log append-only de acciones administrativas. Los campos \texttt{actorEmail} y
\texttt{actorRole} están desnormalizados intencionalmente: reflejan la identidad del actor
en el momento del evento, no su estado actual. Cada registro incluye \texttt{beforeJson} y
\texttt{afterJson}, que permiten reconstruir exactamente qué cambió en cada acción. El
campo \texttt{actorUserId} es nullable con \texttt{onDelete: SetNull}: si el actor es
eliminado, el log persiste con los campos desnormalizados que conservan la identidad
histórica del evento.

\subsection{Relaciones principales}

Las relaciones entre entidades se pueden ver en detalle en el diagrama ER de la
figura~\ref{fig:diagrama-er}. Los aspectos que merecen atención explícita son los
siguientes:

\begin{itemize}
  \item \textbf{Relación \texttt{PlanVersion} $\longrightarrow$ \texttt{Plan} opcional:}
        el campo \texttt{planId} es nullable para soportar versiones de carrera previas a la
        migración del sistema que no tienen un \texttt{Plan} asociado en la nueva tabla
        unificada.

  \item \textbf{\texttt{onDelete: Cascade} en datos de usuario:} \texttt{UserPlanProgress},
        \texttt{ScheduleBlock}, \texttt{UserRecentPlan} y \texttt{UserPreference} se
        eliminan en cascada al borrar el \texttt{User} o la \texttt{PlanVersion}
        padre. Esto evita datos huérfanos sin requerir lógica explícita en la aplicación.

  \item \textbf{FK nula en \texttt{AuditLog}:} \texttt{actorUserId} es nullable con
        \texttt{onDelete: SetNull}. Si el actor es eliminado, el log persiste con los campos
        \texttt{actorEmail} y \texttt{actorRole} desnormalizados, que conservan la identidad
        histórica del evento.
\end{itemize}

\begin{figure}[htbp]
  \centering

  \begin{subfigure}[t]{0.65\textwidth}
    \centering
    \includegraphics[width=\textwidth]{figuras/USER-simplificada.png}
    \caption{Entidades centradas en \texttt{User}.}
    \label{fig:er-user}
  \end{subfigure}

  \vspace{1em}

  \begin{subfigure}[t]{0.65\textwidth}
    \centering
    \includegraphics[width=\textwidth]{figuras/Carrera-simplificada.png}
    \caption{Entidades centradas en \texttt{Carrera}.}
    \label{fig:er-carrera}
  \end{subfigure}

  \vspace{1em}

  \begin{subfigure}[t]{0.65\textwidth}
    \centering
    \includegraphics[width=\textwidth]{figuras/relacion.png}
    \caption{Relación entre \texttt{User} y \texttt{Carrera}.}
    \label{fig:er-relacion}
  \end{subfigure}

  \caption{Diagrama entidad-relación del modelo de datos, organizado en tres bloques.}
  \label{fig:diagrama-er}
\end{figure}

\section{Decisiones de diseño}
\label{sec:decisiones-disenio}

Esta sección justifica las decisiones de diseño más relevantes del sistema, tanto a nivel
de modelo de datos como de selección de herramientas.

\subsection{Modelo de datos}

\begin{itemize}
  \item \textbf{Desnormalización intencional en \texttt{AuditLog}:} los campos
        \texttt{actorEmail} y \texttt{actorRole} se copian literalmente en el momento del
        evento en lugar de referenciarse mediante una clave foránea al usuario. Esta
        decisión protege la integridad histórica del log: si un usuario cambia de rol o
        es eliminado del sistema, el registro de auditoría debe reflejar quién era el actor
        y qué rol tenía en el momento exacto en que ocurrió la acción.

  \item \textbf{Enums \texttt{PlanEstado} y \texttt{PlanFuente} como reemplazo de tablas
        separadas:} la arquitectura anterior usaba tres tablas distintas
        (\texttt{PlanBorrador}, \texttt{PlanPendiente}, \texttt{PlanPublicado}), lo que
        generaba campos inconsistentes entre sí y complicaba las consultas transversales.
        Los enums centralizan esa lógica en una única tabla \texttt{Plan}, simplificando
        migraciones, consultas y el panel de administración.

  \item \textbf{Snapshot inmutable en \texttt{ProgressShare}:} al compartir el avance, el
        sistema copia físicamente el \texttt{stateJson} en ese momento. El enlace siempre
        muestra el estado en el instante en que se compartió, independientemente de cambios
        posteriores del usuario.

  \item \textbf{Progreso del usuario como documento serializado:} el estado de avance
        se almacena como un único objeto JSON (\texttt{stateJson}) en lugar de una tabla
        relacional con una fila por materia. El cliente siempre carga el progreso completo
        de un plan, interactúa localmente y envía el mapa completo al sincronizar: no
        existen consultas del tipo \emph{«dame el estado de la materia X para el usuario Y»}.
        Una tabla normalizada introduciría un \texttt{DELETE} más $N$ inserciones por
        sincronización y reconstruiría en memoria el mismo objeto que ya existe serializado,
        sin beneficio funcional. La estrategia de sincronización LWW y la sanitización del
        payload se describen en la sección~\ref{sec:modulo-progreso}.
\end{itemize}

\subsection{Selección del modelo de IA}
\label{sec:modelos-ia-descartados}

Durante el desarrollo se evaluaron tres alternativas a Gemini para la generación de JSONs
desde PDFs:

\begin{itemize}
  \item \textbf{Grok (xAI):} al momento de evaluación no ofrecía soporte estable para
        \texttt{inlineData} de PDFs y requería extracción de texto previa, lo que contradice
        la estrategia de visión nativa del proyecto.

  \item \textbf{Claude API (Anthropic):} con capacidades de visión comparables a Gemini y
        mejor desempeño en seguimiento de instrucciones complejas. Se descartó por costo:
        el volumen de tokens por PDF resulta significativamente más caro que en Gemini,
        especialmente considerando el plan gratuito disponible durante la etapa de
        desarrollo.

  \item \textbf{Ollama (modelos locales):} permite correr modelos open-source sin costo por
        token, pero los modelos con capacidades de visión suficientes para PDFs complejos
        (LLaVA, Qwen-VL) requieren GPU dedicada. La calidad de extracción en documentos con
        tablas y layouts complejos fue inferior a la de Gemini en las pruebas realizadas.
\end{itemize}

% ============================================================
%  CAPÍTULO 4 - IMPLEMENTACIÓN
% ============================================================
\chapter{Implementación}
\label{cap:implementacion}

\section{Estructura del proyecto}

El repositorio está organizado en tres módulos que se despliegan en servicios separados:
\texttt{frontend/} contiene la aplicación Next.js (App Router, API Routes, componentes de
UI y lógica de plan); \texttt{parser/} contiene el motor de parsing Python con FastAPI,
sus patrones, clasificador y tests por carrera; y \texttt{scripts/} contiene las
utilidades de evaluación, incluyendo \texttt{comparar\_json.py}. Las secciones siguientes
detallan la implementación de los tres módulos principales del sistema.

% ============================================================
%  MÓDULO 1 — PIPELINE GEMINI (VISIÓN NATIVA)
% ============================================================
\section{Módulo: Pipeline Gemini con visión nativa}
\label{sec:modulo-gemini}

\subsection{Flujo de una solicitud}

La figura~\ref{fig:pipelines-comparacion} (al final de la sección~\ref{sec:modulo-parser})
ilustra ambos pipelines en paralelo para facilitar su comparación.

El proceso se inicia en la interfaz de administración, donde el archivo se envía mediante
una solicitud \texttt{POST} hacia la API Route \texttt{/api/admin/planes/parsear} de
Next.js. Debido a las limitaciones de tiempo de ejecución de las funciones
\textit{serverless} en Vercel (\texttt{maxDuration: 120s}), el sistema implementa una
bifurcación de entorno: en producción, la solicitud se delega a un servidor persistente en
Render; en desarrollo local, se procesa directamente desde el runtime de Next.js.

El núcleo del pipeline reside en la integración con la API de Google Generative AI, donde
el PDF se transmite como \texttt{inlineData} en base64 con MIME type
\texttt{application/pdf}. Esta técnica permite aprovechar la visión nativa del modelo
Gemini~2.5~Flash para interpretar la jerarquía visual del documento --- tablas, encabezados
y saltos de página --- sin requerir una capa previa de extracción de texto plano. Para
garantizar resultados predecibles y repetibles, el modelo opera con temperatura~0. La
temperatura es un parámetro que controla el grado de aleatoriedad en la generación de
texto: con temperatura alta el modelo explora opciones menos probables y produce respuestas
más variadas y creativas; con temperatura~0 selecciona siempre el token de mayor
probabilidad en cada paso, lo que produce una salida determinista. Para una tarea de
extracción estructurada como esta --- donde el output correcto es único y está definido por
un schema estricto --- la creatividad es indeseable: dos ejecuciones sobre el mismo PDF
deben producir el mismo JSON. Temperatura~0 elimina esa variabilidad y hace los resultados
reproducibles y comparables entre ejecuciones.

La respuesta cruda del modelo es procesada por la función \texttt{extraerJSON}, que limpia
posibles bloques de Markdown y aplica tres intentos de parseo en cascada para asegurar la
validez del formato. El JSON resultante se muestra en el panel de comparación pero
\textbf{no se persiste automáticamente}: el admin decide explícitamente si guardarlo como
borrador o publicarlo directamente (véase la subsección~\ref{sec:guardado-resultado}).

Este flujo ocurre en paralelo al del parser Python (sección~\ref{sec:modulo-parser}): el
panel admin lanza ambas solicitudes simultáneamente al recibir el PDF, y presenta los dos
JSONs resultantes en una vista de comparación side-by-side que se describe en la
sección~\ref{sec:integracion-flujo}.

\subsection{Estructura del prompt (v33)}

El prompt del sistema (constante \texttt{DEFAULT\_SYSTEM\_PROMPT} en
\texttt{frontend/lib/ai/prompt.ts}) tiene la siguiente estructura:

\begin{enumerate}
  \item \textbf{Rol del modelo:} se define como ``motor de extracción determinista
        de planes de estudio''.
  \item \textbf{Instrucción crítica de continuidad entre páginas:} dado que los
        PDFs de la UNS frecuentemente cortan la tabla de correlativas de una materia
        entre dos páginas, la primera instrucción del prompt explica que las filas
        que comienzan una página sin nombre de materia pertenecen \emph{siempre} a la
        última materia de la página anterior. Esta instrucción apareció en versiones
        tempranas del prompt y se reforzó luego de detectar que era uno de los errores
        más frecuentes del modelo.
  \item \textbf{Schema de salida:} el JSON target completo con tipos y valores
        válidos para cada campo.
  \item \textbf{Detalles de campo:} sección por sección, las reglas de normalización
        (e.g.\ los valores de \texttt{para\_cursar} y \texttt{para\_rendir} deben
        ser siempre minúscula: \texttt{"cursada"} o \texttt{"aprobada"}, aunque el PDF
        los muestre capitalizados).
  \item \textbf{Patrones especiales (CRITICAL):} dos patrones que requieren
        generación condicional de entradas en dos arreglos simultáneos:
        \begin{itemize}
          \item \textbf{POSITION A / POSITION B} para agrupadores de optativas
                (\texttt{G\#\#\#\#}): un grupo en POSITION~A debe generar una entrada
                en \texttt{materias[]} \emph{y} otra en \texttt{agrupadores[]};
                en POSITION~B solo genera entrada en \texttt{agrupadores[]}.
          \item \textbf{Grupos de idioma} (\texttt{I\#\#\#\#}): el prompt incluye
                una instrucción explícita para que el modelo distinga la letra
                \texttt{I} mayúscula del dígito \texttt{1} --- un error sistemático
                detectado durante la evaluación. Se incluyó además una instrucción
                de validación cruzada al final: antes de devolver el JSON, el modelo
                debe escanear todos los \texttt{correlativas} y reemplazar claves
                numéricas como \texttt{"12201"} por su equivalente correcto
                \texttt{"I2201"}.
        \end{itemize}
  \item \textbf{Reglas globales y validación:} lista de invariantes que debe
        cumplir el JSON (arrays no vacíos, objetos en lugar de null, sin duplicados,
        etc.).
\end{enumerate}

El prompt es editable en tiempo de ejecución desde la tabla \texttt{AdminConfig}
de la base de datos. Si existe una fila con \texttt{id = "singleton"}, se usa ese
valor; si no, se cae al texto compilado en \texttt{prompt.ts}. Esto permite iterar
el prompt sin necesidad de redeploy.

\subsection{Extracción del JSON de la respuesta}

El modelo puede devolver el JSON directamente, envuelto en un bloque markdown
(\texttt{```json ... ```}), o con texto adicional antes o después. La función
\texttt{extraerJSON} implementa tres intentos en cascada: parseo directo,
extracción del bloque markdown, y extracción por delimitadores de llaves.
Si los tres fallan, devuelve el error con un preview de los primeros 500
caracteres de la respuesta.

\subsection{Guardado del resultado}
\label{sec:guardado-resultado}

Una vez extraído el JSON, el resultado se muestra en el panel de comparación pero
\textbf{no se persiste automáticamente}. El admin decide explícitamente qué hacer
con el resultado mediante dos acciones:

\begin{itemize}
  \item \textbf{Guardar borrador:} persiste el plan con \texttt{estado = BORRADOR}
        para revisión posterior, sin publicarlo todavía.
  \item \textbf{Publicar:} persiste el plan directamente con
        \texttt{estado = PUBLICADO}, haciéndolo visible para los usuarios de
        inmediato.
\end{itemize}

En ambos casos el plan queda registrado en la tabla \texttt{Plan} con
\texttt{fuente = GEMINI}. Esta decisión mantiene la base de datos limpia: solo
se persiste lo que el admin consideró suficientemente válido para conservar.

% ============================================================
%  MÓDULO 2 — PARSER LOCAL EN PYTHON
% ============================================================
\section{Módulo: Parser local en Python}
\label{sec:modulo-parser}

\subsection{Responsabilidad y posición en el sistema}

El parser local cumple el rol de generador del \textit{ground truth}: produce el
JSON de referencia contra el cual se evalúa la calidad del output de Gemini. No
es un adapter que corrige el output del modelo; es una fuente de verdad
independiente basada en reglas explícitas y verificables.

\subsection{Flujo de una solicitud}

El pipeline del parser local se ilustra en la figura~\ref{fig:pipeline-python},
presentada al final de esta sección junto al pipeline Gemini
(figura~\ref{fig:pipeline-gemini}) para facilitar la comparación de ambos flujos.

El proceso se inicia en \texttt{CargarPlanTab.tsx} cuando el administrador
selecciona el modo ``Parser local''. La acción desencadena la cadena
\texttt{handleFile} $\rightarrow$ \texttt{chequearYParsear("parser")} $\rightarrow$
\texttt{ejecutarParseo} $\rightarrow$ \texttt{parsearSSE}, que envía el PDF
mediante \texttt{POST} a la API Route \texttt{/api/admin/planes/parsear-local}.

Al igual que el pipeline Gemini, la route bifurca según el entorno:

\begin{itemize}
  \item \textbf{Producción (Render):} si la variable de entorno
        \texttt{PARSER\_API\_URL} está definida, el PDF se reenvía al servicio
        FastAPI en Render (\texttt{POST \$\{PARSER\_API\_URL\}/parse}). La respuesta
        llega como JSON directo con \texttt{\{type: "done", data\}}.

  \item \textbf{Desarrollo local:} si \texttt{PARSER\_API\_URL} no está definida,
        la route escribe el PDF a \texttt{/tmp} y ejecuta un subprocess:
        \texttt{python3 -m core.parser <pdfPath> <outPath>}. La respuesta llega
        como stream de eventos SSE (\texttt{guardando} $\rightarrow$
        \texttt{parseando} $\rightarrow$ \texttt{leyendo} $\rightarrow$
        \texttt{done}).
\end{itemize}

En ambos casos, el JSON resultante se muestra en el panel de comparación pero
no se persiste automáticamente. Al igual que en el pipeline Gemini, el admin
decide explícitamente qué hacer con el resultado: guardarlo como borrador
(\texttt{estado = BORRADOR}) para revisión posterior, o publicarlo directamente
(\texttt{estado = PUBLICADO}). En ambos casos el plan queda registrado en Neon
con \texttt{fuente = PARSER}.

\subsection{Motor de extracción: tres pasos}

Independientemente del entorno, el núcleo del parsing es siempre
\texttt{parsear\_plan\_pdf()} en \texttt{cli.py}, que ejecuta tres pasos en
secuencia:

\begin{enumerate}
  \item \textbf{Extracción de texto plano} (\texttt{extractor.py}):
        \texttt{pdfplumber} abre el PDF y llama a \texttt{page.extract\_text()}
        en cada página, concatenando los resultados con saltos de línea. El texto
        resultante preserva el orden visual de las líneas tal como las organiza el
        motor de layout de \texttt{pdfplumber}.

  \item \textbf{Limpieza y recomposición} (\texttt{cleaner.py}): elimina líneas
        de decoración (encabezados de tabla como \texttt{"Para cursar Para rendir"},
        pie de página como \texttt{"Ir a Moodle UNS"}, líneas de navegación como
        \texttt{"Volver"}) usando cadenas exactas y prefijos conocidos. Luego
        \texttt{recomponer\_lineas\_partidas} fusiona líneas que \texttt{pdfplumber}
        partió incorrectamente: correlativas sin su segundo estado, nombres de
        materias cortados en dos líneas, y prosa de requisitos especiales partida
        en el salto de página.

  \item \textbf{Clasificación y construcción del JSON}
        (\texttt{parser\_plan.py}, \texttt{classifiers.py}): cada línea pasa por
        \texttt{clasificar\_linea()}, que la categoriza con los patrones regex de
        \texttt{patterns.py}. La tabla~\ref{tab:patrones} muestra los tipos y
        ejemplos representativos. Según el tipo, el parser actualiza el estado de
        contexto (\texttt{año\_actual}, \texttt{cuatrimestre\_actual},
        \texttt{orientacion\_actual}) y construye los arreglos \texttt{materias[]}
        y \texttt{agrupadores[]}. Al terminar el recorrido línea a línea, una
        segunda pasada resuelve las orientaciones por materia.
\end{enumerate}

\begin{table}[htbp]
\centering
\small
\begin{tabular}{p{3.5cm} p{4.5cm} p{5cm}}
\toprule
\textbf{Tipo} & \textbf{Ejemplo} & \textbf{Patrón (simplificado)} \\
\midrule
\texttt{anio}
  & \texttt{PRIMER AÑO}
  & \texttt{/\^{}(PRIMER|...|SEXTO)\textbackslash{}s+AÑO\$/i} \\
\addlinespace
\texttt{cuatrimestre}
  & \texttt{Primer Cuatrimestre}
  & \texttt{/\^{}(Anual|Primer|Segundo Cuatrimestre)\$/i} \\
\addlinespace
\texttt{materia}
  & \texttt{3942 ALGEBRA 96hs.}
  & \texttt{/\^{}([A-Z]?\textbackslash{}d\{3,\})\textbackslash{}s+(.+?)\$/i} \\
\addlinespace
\texttt{correlativa}
  & \texttt{2405 Aprobada Aprobada}
  & \texttt{/([A-Z]?\textbackslash{}d\{3,\})\textbackslash{}s+(Aprobada|Cursada)\{2\}/i} \\
\addlinespace
\texttt{seccion\_optativas}
  & \texttt{MATERIAS OPTATIVAS}
  & \texttt{/\^{}MATERIAS\textbackslash{}s+OPTATIVAS\$/i} \\
\addlinespace
\texttt{seccion\_orientacion}
  & \texttt{ORIENTACIÓN Hidráulica}
  & \texttt{/\^{}ORIENTACI[OÓ]N\textbackslash{}s+(.+)\$/i} \\
\addlinespace
\texttt{grupo}
  & \texttt{G0624 Nombre grupo}
  & \texttt{/\^{}([A-Z]\textbackslash{}d\{4,\})\textbackslash{}s*(.*)\$/i}
    (solo en sección optativas) \\
\bottomrule
\end{tabular}
\caption{Tipos de línea del clasificador del parser local con ejemplos reales.}
\label{tab:patrones}
\end{table}

\subsection{Validación de schema y diferencia entre entornos}

En desarrollo local, \texttt{contract\_validator.py} valida el JSON generado
contra el schema del sistema antes de escribirlo a disco. Si la validación falla,
el proceso termina con código~1 y la route Next.js lo reporta como un evento de
error SSE visible en el panel admin. En producción (Render), esta validación no
corre: \texttt{parsear\_plan\_pdf} devuelve el diccionario directamente y Render
responde con \texttt{\{type: "done", data\}}. La validación queda delegada al
panel de comparación, donde el admin puede detectar anomalías antes de publicar.

\subsection{Heurísticas de contexto}

La mayor parte de la lógica del parser no está en los patrones sino en cómo el
estado de contexto se actualiza y propaga:

\begin{itemize}
  \item \textbf{Estado de sección:} \texttt{seccion\_actual} puede ser
        \texttt{"normal"}, \texttt{"optativas"}, \texttt{"idiomas"} o
        \texttt{"seminarios"}. Este estado controla qué tipos de línea son válidos
        en cada momento: un patrón de grupo (\texttt{G\#\#\#\#}) solo se
        interpreta como encabezado de grupo optativo dentro de la sección
        \texttt{"optativas"}.

  \item \textbf{Colisiones de ID:} algunos IDs aparecen primero como materia en
        la sección normal y luego como encabezado de agrupador en la sección
        optativas (o viceversa). El parser detecta estas colisiones: cuando un ID
        ya registrado en \texttt{materias\_index} aparece como encabezado de grupo,
        transfiere su ubicación al agrupador y elimina la entrada duplicada de
        \texttt{materias[]}. La excepción son las colisiones de tipo
        \texttt{idioma\_grupo}: el ID debe permanecer en \texttt{materias[]} con
        \texttt{tipo: "materia"}, ya que la regla POSITION~A así lo requiere.

  \item \textbf{Orientaciones múltiples:} cuando el PDF divide el plan en bloques
        de orientación, el parser rastrea en qué orientación apareció cada materia
        usando dos estructuras separadas para materias obligatorias y optativas. Al
        finalizar el recorrido, las materias presentes en todos los bloques se
        marcan como comunes; las que solo aparecen en algunos reciben el campo
        \texttt{orientacion} o \texttt{orientaciones} según corresponda.

  \item \textbf{Correlativas en prosa:} las líneas que no pueden tipificarse como
        correlativa estructurada (\texttt{ID Estado Estado}) pero contienen texto de
        requisito son procesadas por \texttt{correlativa\_prosa.py}, que intenta
        inferir el \texttt{requisito\_especial[]} correspondiente o, como fallback,
        genera correlativas aproximadas con un warning.
\end{itemize}

\begin{figure}[H]
  \centering
  \begin{subfigure}[t]{0.47\textwidth}
    \centering
    \includegraphics[width=\textwidth]{figuras/pipeline-gemini.png}
    \caption{Pipeline Gemini (visión nativa).}
    \label{fig:pipeline-gemini}
  \end{subfigure}
  \hfill
  \begin{subfigure}[t]{0.47\textwidth}
    \centering
    \includegraphics[width=\textwidth]{figuras/pipeline-python.png}
    \caption{Pipeline parser local (Python).}
    \label{fig:pipeline-python}
  \end{subfigure}
  \caption{Pipelines de extracción del sistema. Ambos se ejecutan en paralelo al subir
           un PDF desde el panel admin y sus resultados se comparan side-by-side antes
           de publicar.}
  \label{fig:pipelines-comparacion}
\end{figure}

% ============================================================
%  MÓDULO 3 — SINCRONIZACIÓN DE PROGRESO
% ============================================================
\section{Módulo: Sincronización de progreso del usuario}
\label{sec:modulo-progreso}

\subsection{Modelo de almacenamiento}

El progreso del usuario en un plan se almacena como un único objeto JSON
serializado en el campo \texttt{stateJson} de la tabla \texttt{UserPlanProgress}.
Su estructura es un mapa plano de ID de materia a estado:

\begin{lstlisting}[language=JavaScript,
                   caption={Estructura del \texttt{stateJson} en \texttt{UserPlanProgress}.}]
{
  "5175": "aprobada",
  "3942": "cursada",
  "4102": "no_cursada",
  ...
}
\end{lstlisting}

Esta decisión está justificada por el patrón de acceso real del sistema: el
cliente siempre carga el progreso completo de un plan al abrir la vista, el
usuario interactúa localmente con los estados de cada materia, y al sincronizar
envía el mapa completo de vuelta al servidor. No existen consultas del tipo
\emph{«dame el estado de la materia X para el usuario Y»}: la unidad mínima de
lectura y escritura es siempre el progreso completo del plan.

Una tabla normalizada (\texttt{EstadoMateria} con una fila por materia por usuario)
introduciría complejidad operativa sin beneficio funcional: cada sincronización
requeriría un \texttt{DELETE} seguido de $N$ inserciones (o un upsert por materia),
y cada lectura reconstruiría en memoria el mismo objeto que ya existe serializado.

\subsection{Protocolo de sincronización: Last-Write-Wins}

El cliente mantiene una copia local del progreso en el almacenamiento del
navegador, con un timestamp \texttt{updatedAt}. Al sincronizar, el endpoint
\texttt{POST /api/progreso/sync} aplica una política
\textit{Last-Write-Wins} (LWW) basada en ese timestamp:

\begin{lstlisting}[language=JavaScript,
                   caption={Resolución LWW en \texttt{progressSync.ts}.}]
export function resolveProgressSnapshotLww(params: {
  local: ProgressSnapshot;
  remote: ProgressSnapshot;
}): ProgressResolution {
  const localMs  = toMillis(params.local.updatedAt);
  const remoteMs = toMillis(params.remote.updatedAt);

  if (localMs >= remoteMs) {
    return { source: "local",  snapshot: params.local };
  }
  return { source: "remote", snapshot: params.remote };
}
\end{lstlisting}

Si el snapshot local es más reciente (o igual), el servidor persiste el estado
local mediante un \texttt{upsert} sobre la clave compuesta
\texttt{(userId, planVersionId)}. Si el snapshot remoto es más reciente
(por ejemplo, el usuario actualizó su progreso desde otro dispositivo), el
servidor devuelve el estado remoto y el cliente lo adopta.

\subsection{Sanitización del payload}

Antes de persistir, la función \texttt{sanitizeProgressState} valida el payload
descartando cualquier entrada cuya clave sea vacía, cuyo valor no sea un string,
o cuyo valor no sea uno de los tres estados válidos (\texttt{"no\_cursada"},
\texttt{"cursada"}, \texttt{"aprobada"}). Esto garantiza que nunca lleguen a la
base de datos estados inválidos provenientes de versiones antiguas del cliente o
de manipulación del payload.

\subsection{Trazabilidad del upsert}

Cuando el servidor persiste un estado local (es decir, cuando el cliente tenía
datos más recientes que la base de datos), el endpoint registra un evento
\texttt{PROGRESS\_MIGRATED} en la tabla \texttt{AuditLog}, capturando el estado
anterior y posterior del progreso. Esta decisión permite al equipo de desarrollo
detectar situaciones inusuales (por ejemplo, pérdida masiva de progreso) sin
exponer los datos del usuario al cliente.

\section{Integración y flujo general}
\label{sec:integracion-flujo}

El flujo completo desde la carga de un PDF hasta la publicación del plan
es el siguiente:

\begin{enumerate}
  \item El administrador sube un PDF desde el panel \texttt{/admin}.
  \item El frontend realiza dos requests en paralelo: uno a
        \texttt{/api/admin/planes/parsear} (Gemini) y otro a
        \texttt{/api/admin/planes/parsear-local} (parser Python en Render).
  \item Ambos resultados se muestran en una vista de comparación side-by-side.
        El admin puede inspeccionar las diferencias campo a campo mediante el
        componente \texttt{MergeInteractivo}.
  \item El admin elige qué JSON publicar (o fusiona manualmente campos de ambos),
        y confirma desde el drawer de publicación (\texttt{GuardarPlanDrawer}).
  \item El plan queda registrado en la tabla \texttt{Plan} con
        \texttt{estado = PUBLICADO} y es inmediatamente visible para los usuarios
        en \texttt{/planes/[carrera]}.
\end{enumerate}

% ============================================================
%  CAPÍTULO 5 - EVALUACIÓN Y RESULTADOS
% ============================================================
\chapter{Evaluación y Resultados}
\label{cap:evaluacion}

\section{Metodología de evaluación}

[Describir cómo se evaluó el sistema: pruebas funcionales, pruebas con usuarios,
benchmarks de rendimiento, etc.]

\section{Resultados obtenidos}

[Presentar los resultados. Podés usar tablas, gráficos, capturas de pantalla, etc.]

\section{Problemas encontrados al usar Gemini}
\label{sec:problemas-gemini}

La integración con la API de Gemini presentó varios problemas concretos que condicionaron
las decisiones de diseño del pipeline y que se detectaron durante la evaluación comparativa
contra el ground truth del parser local.

\textbf{Restricciones de la API gratuita.} Durante el desarrollo se usó el plan gratuito
de Google AI Studio, que impone límites de tasa por minuto, por hora y por día: para
\texttt{gemini-2.5-flash}, el plan gratuito permite 10~requests por minuto y 500 por día.
Estos límites son suficientes para uso administrativo ocasional (subir un PDF y esperar el
JSON), pero no para procesamiento batch de múltiples carreras en paralelo. El panel admin
incluye manejo explícito de estos límites y muestra el estado de rate limit al usuario
cuando se alcanza.

\textbf{Calidad variable entre modelos.} \texttt{gemini-2.5-flash} resultó el mejor
balance entre calidad de extracción y velocidad de respuesta. \texttt{gemini-2.5-pro}
produce outputs de mayor calidad pero es significativamente más lento y tiene límites de
rate más estrictos (5~requests/minuto en el plan gratuito). \texttt{gemma-4-26b-a4b-it},
el modelo open-weight de Google disponible vía API, resultó impráctico para el flujo del
panel admin: tiempos de respuesta de varios minutos por PDF y calidad de extracción
inferior en documentos con tablas complejas.

\textbf{Problema de salto de página.} Gemini tiene dificultades para detectar correctamente
las correlativas de materias cuya entrada en el PDF está dividida entre dos páginas. Cuando
el nombre de una materia aparece al final de una página y sus correlativas al inicio de la
siguiente, el modelo frecuentemente omite las correlativas o las asocia a la materia
incorrecta. Este es uno de los errores más frecuentes detectados en la comparación con el
parser local.

\textbf{Confusión entre \texttt{I} y \texttt{1} en IDs.} Los identificadores de grupos de
idioma comienzan con la letra \texttt{I} mayúscula (ej: \texttt{I0001}, \texttt{I0002}).
Gemini confunde sistemáticamente esta letra con el número \texttt{1}, generando IDs
inválidos como \texttt{10001}. Este error es consistente y reproducible, y requirió una
instrucción explícita en el prompt del sistema para mitigarlo, aunque no se eliminó
completamente.

\section{Discusión}

[Analizar los resultados: qué funcionó bien, qué limitaciones se encontraron, cómo se
compara con el estado del arte.]

\subsection{Escalabilidad de la base de datos y límites del plan gratuito de Neon}

El diseño del schema de la base de datos es adecuado para soportar múltiples usuarios
concurrentes en el volumen esperado para la UNS. Las decisiones de diseño más relevantes
para la concurrencia son las siguientes:

\begin{itemize}
  \item Todas las tablas de usuario tienen índices sobre \texttt{(userId, createdAt)},
        lo que garantiza que las consultas por usuario sean O(log n) sin necesidad de
        full scans aunque la base de datos crezca.
  \item \texttt{UserPlanProgress} almacena una única fila por \texttt{(userId,
        planVersionId)}: no hay contención entre usuarios porque cada uno escribe
        exclusivamente sobre su propia fila.
  \item \texttt{AuditLog} es una tabla append-only: las inserciones no generan locks
        sobre filas existentes y múltiples usuarios pueden escribir en paralelo sin
        interferencia.
  \item El connection pooling vía PgBouncer (configurado en Neon) permite hasta 10.000
        conexiones concurrentes independientemente del tamaño del compute, absorbiendo
        picos de tráfico sin agotar el pool de conexiones de PostgreSQL.
\end{itemize}

\textbf{Límites del plan gratuito de Neon.} El sistema utiliza el plan gratuito de Neon
durante su etapa de desarrollo y despliegue inicial. Este plan impone tres límites que
son relevantes para la escalabilidad:

\begin{itemize}
  \item \textbf{Storage: 0,5 GB.} Si se supera este límite, el compute queda suspendido
        hasta el siguiente ciclo de facturación. Para el volumen actual del sistema
        (decenas de carreras con JSONs de plan de pocos kilobytes cada uno, más filas de
        progreso de usuarios) este límite es holgado, pero podría convertirse en un
        cuello de botella si el sistema se generaliza a múltiples universidades con
        grandes volúmenes de datos de actividad.
  \item \textbf{Compute hours: 100 CU-hours por mes.} Equivale aproximadamente a una
        instancia de 0,25 CU corriendo durante 400 horas continuas. Para uso
        administrativo esporádico este presupuesto es más que suficiente, pero con
        cientos de usuarios activos simultáneos el compute se agotaría antes de fin
        de mes, dejando la aplicación sin base de datos hasta el siguiente ciclo.
  \item \textbf{Scale-to-zero automático.} La instancia de Neon entra en suspensión
        tras 5 minutos de inactividad (no configurable en el plan gratuito). El primer
        request tras el período de inactividad experimenta un cold start de 1--2
        segundos mientras el compute vuelve a levantarse. Para la aplicación web esto
        es perceptible como una latencia inicial inusual pero no constituye un error.
\end{itemize}

En conclusión, el plan gratuito de Neon es adecuado para el uso actual del sistema
(panel admin con tráfico bajo y usuarios universitarios de una sola institución), pero
no es suficiente si la aplicación escala a miles de usuarios activos simultáneos o se
generaliza a múltiples universidades. En ese escenario sería necesario migrar al plan
pago de Neon (que elimina los límites de compute y aumenta el storage) o evaluar otras
opciones de PostgreSQL gestionado como Supabase o Railway.

\textbf{Problemas de performance a futuro.} Más allá de los límites de Neon, existen
dos aspectos del diseño actual que podrían generar problemas de rendimiento a escala:

\begin{itemize}
  \item \textbf{\texttt{planJson} como texto sin comprimir:} el JSON completo de cada
        carrera se almacena como string en la tabla \texttt{Plan}. Con planes de 200+
        materias y muchos usuarios leyendo el mismo plan en paralelo, el volumen de
        datos transferidos por query aumenta. Una optimización posible es mover el JSON
        a almacenamiento de objetos (S3, R2 o similar) y guardar solo la URL en la base
        de datos, reduciendo la carga sobre PostgreSQL.
  \item \textbf{Eliminación de \texttt{UserActivity}:} una versión anterior del schema
        incluía una tabla \texttt{UserActivity} que registraba en forma append-only todos
        los eventos del usuario (materias marcadas, cambios de carrera activa, aperturas
        de plan). Se eliminó porque su crecimiento era ilimitado sin ninguna política de
        retención definida, y las funcionalidades que habilitaba (feed de actividad en el
        perfil, estadísticas de planes más consultados en el dashboard) eran prescindibles
        frente al costo de mantener una tabla que crece indefinidamente. El único dato de
        trazabilidad que se conserva es \texttt{AuditLog}, que registra exclusivamente
        acciones administrativas con alcance acotado y un volumen de escritura bajo.
\end{itemize}

\subsection{Extensión multi-institucional: incorporación de otras universidades}

Una extensión natural del sistema es soportar planes de estudio de universidades distintas
a la UNS. El schema actual está diseñado alrededor de una única institución implícita: la
jerarquía \texttt{Departamento} $\rightarrow$ \texttt{Carrera} $\rightarrow$
\texttt{PlanVersion} no incluye un nivel de universidad porque el sistema nació como
una herramienta específica para la UNS.

Agregar soporte multi-institucional requeriría una modificación estructural mínima: crear
una tabla \texttt{Universidad} e incorporar una clave foránea \texttt{universidadId} en
\texttt{Departamento}, extendiendo la jerarquía a:

\begin{center}
\texttt{Universidad} $\longrightarrow$ \texttt{Departamento} $\longrightarrow$
\texttt{Carrera} $\longrightarrow$ \texttt{PlanVersion} $\longrightarrow$ \texttt{Plan}
\end{center}

Todo lo demás del schema --- progreso de usuarios, bloques de horario, auditoría ---
permanece sin cambios, ya que esas tablas referencian \texttt{Carrera} o
\texttt{PlanVersion} directamente y son independientes de la institución.

El aspecto más delicado de esta extensión no es estructural sino de unicidad. El campo
\texttt{id} de \texttt{Carrera} es hoy un slug único globalmente
(ej: \texttt{ingenieria\_civil}). Con múltiples universidades pueden existir carreras con
el mismo nombre en distintas instituciones, lo que generaría colisiones de slug. La
solución es cambiar la unicidad de global a institucional, por ejemplo mediante un slug
compuesto (\texttt{uns\_\_ingenieria\_civil}, \texttt{utn\_\_ingenieria\_civil}) o
modificando la restricción de unicidad a \texttt{@@unique([universidadId, slug])} y
actualizando las rutas de la aplicación para incluir el identificador de universidad
(\texttt{/planes/[universidad]/[carrera]}).

En síntesis, la base de datos soporta esta extensión sin rediseño profundo: es una
migración incremental que agrega un nivel en la jerarquía existente, con el único cambio
no trivial en la estrategia de slugs de carrera.

\subsection{Versionado de planes y migración entre versiones}

El schema del sistema incluye la tabla \texttt{PlanVersion} precisamente para soportar
el escenario en que una universidad actualiza su plan de estudios y coexisten múltiples
versiones activas. Un estudiante que inició su carrera con el plan 2012 (\texttt{v1}) puede
seguir trabajando con esa versión aunque la universidad publique el plan 2027 (\texttt{v2}),
ya que su progreso está anclado a una \texttt{PlanVersion} específica mediante la clave
\texttt{(userId, planVersionId)} en \texttt{UserPlanProgress}.

Sin embargo, la implementación actual de este mecanismo está incompleta. La opción
``Guardar como nueva versión'' existe en la interfaz del panel admin pero está
deshabilitada: el backend siempre muta el \texttt{Plan} existente en lugar de crear una
nueva \texttt{PlanVersion} y actualizar \texttt{Carrera.defaultVersionId}. Esta
limitación está identificada como trabajo pendiente en el sistema y no requiere cambios
en el schema de la base de datos para resolverse, solo completar la lógica del backend y
la experiencia de usuario del panel admin.

La base de datos soporta el versionado múltiple sin modificaciones: la restricción
\texttt{@@unique([carreraId, versionId])} en \texttt{PlanVersion} permite que coexistan
\texttt{v1} y \texttt{v2} para la misma carrera, cada una con su propio \texttt{Plan}
publicado y con el progreso de cada usuario preservado de forma independiente.

% ============================================================
%  CAPÍTULO 6 - CONCLUSIÓN
% ============================================================
\chapter{Conclusión}
\label{cap:conclusiones}

Este proyecto logró desarrollar con éxito un sistema integral para la transformación de
planes de estudio universitarios estáticos en herramientas digitales interactivas. La
solución abordó de forma efectiva el problema de la brecha entre el formato PDF y los datos
estructurados mediante un enfoque híbrido que combina la precisión de un parser basado en
reglas con la flexibilidad de la inteligencia artificial generativa.

\section{Aportes del trabajo}

\begin{itemize}
  \item \textbf{Pipeline de extracción híbrido:} Se implementó una arquitectura capaz de
        procesar documentos heterogéneos de la UNS, utilizando un parser local en Python
        para generar el \textit{ground truth} y el modelo Gemini para una conversión
        escalable mediante visión nativa.

  \item \textbf{Esquema de datos robusto:} Se diseñó un schema JSON validado que captura
        la complejidad académica (correlativas, optativas e idiomas) y un modelo de base de
        datos optimizado para la trazabilidad y el rendimiento.

  \item \textbf{Aplicación de seguimiento interactivo:} Se desarrolló una plataforma web
        que permite a los estudiantes de la UNS visualizar dependencias de forma gráfica y
        gestionar su progreso sincronizado en la nube.

  \item \textbf{Métrica de validación de IA:} Se estableció un mecanismo de comparación
        entre el parser y el LLM, permitiendo a los administradores auditar la precisión de
        la IA antes de la publicación oficial de los datos.
\end{itemize}

\section{Lecciones aprendidas}

A lo largo del desarrollo de este proyecto aprendí [reflexión personal sobre el proceso,
tecnologías nuevas, desafíos superados, etc.].

\section{Trabajo futuro}

Las siguientes líneas de trabajo quedaron fuera del alcance de esta versión pero están
identificadas como extensiones naturales del sistema:

\begin{enumerate}
  \item \textbf{Generalización multi-institucional:} Expandir el sistema para soportar
        planes de otras universidades nacionales, aprovechando la configuración de
        \textit{prompts} genéricos ya prevista en el sistema.

  \item \textbf{Analítica avanzada de datos:} Implementar procesos de agregación sobre los
        JSON de progreso para generar estadísticas sobre tasas de recursado y materias
        críticas, sin alterar el modelo de almacenamiento actual.

  \item \textbf{Integración con sistemas de gestión:} Explorar la vinculación mediante API
        con sistemas de autogestión de alumnos (como SIU-Guaraní) para automatizar la
        carga inicial del progreso académico.

  \item \textbf{Historia académica automática:} Permitir que el usuario suba el PDF de su
        historia académica y que el sistema marque automáticamente las materias aprobadas y
        cursadas en el plan, sin necesidad de ingresarlas manualmente una a una.

  \item \textbf{Nota y promedio por materia:} Ampliar el modelo de progreso para que el
        usuario pueda registrar la nota obtenida en cada materia aprobada. Con esa
        información el sistema podría calcular el promedio académico y proyectarlo a futuro.

  \item \textbf{Mejora paralela del prompt y el parser:} El parser local y el prompt de
        Gemini evolucionan de forma independiente hoy. Una línea de mejora es usar los
        errores detectados por el comparador como señal de retroalimentación para iterar
        ambos en paralelo, construyendo un dataset incremental donde cada nuevo plan
        correcto se convierte en un caso de prueba adicional.

  \item \textbf{Integración completa de carreras pendientes:} La UNS tiene decenas de
        carreras aún no incorporadas al sistema. Completar la cobertura requiere procesar
        los PDFs restantes, validar los JSONs generados y publicarlos en el panel admin.

  \item \textbf{Evaluación automática de requisitos especiales:} Implementar la evaluación
        de las condiciones de mínimo de materias aprobadas y año completo aprobado
        directamente sobre el progreso marcado por el usuario, para que el sistema pueda
        bloquear o habilitar materias con requisitos especiales de forma automática.

  \item \textbf{Soporte de cursado paralelo:} Extender el schema JSON y la lógica de
        habilitación para representar correlativas que admiten inscripción simultánea,
        cuando el reglamento de la carrera lo permite.

  \item \textbf{Selector de versiones de plan y migración asistida:} Completar el flujo de
        publicación de nuevas versiones en el panel admin y diseñar la experiencia de
        selección de versión para los usuarios finales, de modo que estudiantes de distintos
        años de ingreso puedan trabajar con el plan que les corresponde. Como extensión de
        esta funcionalidad, implementar una vista de comparación entre versiones que permita
        al usuario visualizar qué materias tienen en común \texttt{v1} y \texttt{v2},
        proyectar su progreso actual sobre el nuevo plan (identificando qué materias ya
        aprobadas son reconocidas en la versión nueva), y tomar una decisión informada sobre
        si conviene o no migrar al plan actualizado.
\end{enumerate}

% ============================================================
%  BIBLIOGRAFÍA
% ============================================================
\newpage
\chapter*{Bibliografía}
\addcontentsline{toc}{chapter}{Bibliografía}

\begin{thebibliography}{99}

\bibitem{ref1}
[Autor/es].
\textit{[Título del libro o artículo]}.
[Editorial / Revista / Sitio web], [Año].
Recuperado de: \url{https://ejemplo.com}

\bibitem{ref2}
[Autor/es].
\textit{[Título]}.
[Fuente], [Año].

\bibitem{ref3}
Universidad Nacional del Sur.
\textit{Planes de carreras}.
Recuperado de: \url{https://www.uns.edu.ar/academicas/carreras/oferta-academica_carreras-grado}

\bibitem{neon-pricing}
Neon.
\textit{Pricing --- Neon Docs}.
2025.
Recuperado de: \url{https://neon.com/pricing}

\bibitem{neon-plans}
Neon.
\textit{Plans --- Neon Docs}.
2025.
Recuperado de: \url{https://neon.com/docs/introduction/plans}

\bibitem{neon-pooling}
Neon.
\textit{Connection pooling --- Neon Docs}.
2025.
Recuperado de: \url{https://neon.com/docs/connect/connection-pooling}

\end{thebibliography}

% ============================================================
%  APÉNDICES
% ============================================================
\appendix

\chapter{Glosario}
\label{ap:glosario}

\begin{description}[style=nextline]
  \item[Agrupador] Entidad del schema JSON que representa un grupo de materias optativas
        o de idioma. Puede aparecer como correlativa de una materia obligatoria
        (POSITION~A) o solo como encabezado de sección (POSITION~B).
  \item[Ground truth] JSON generado por el parser local en Python, usado como
        referencia para evaluar la calidad del output de Gemini.
  \item[inlineData] Mecanismo de la API de Gemini para enviar el contenido de un
        archivo (PDF, imagen) codificado en base64 directamente en el cuerpo del
        request, sin cargarlo previamente a un servicio de almacenamiento.
  \item[LWW (Last-Write-Wins)] Política de resolución de conflictos en la que,
        ante dos versiones en conflicto, se conserva la más reciente según su
        timestamp.
  \item[Slug] Identificador de URL de una carrera, derivado de su nombre en
        minúsculas y sin acentos, con espacios reemplazados por guiones bajos
        (ej: \texttt{ingenieria\_civil}).
  \item[Visión nativa] Capacidad de un modelo de lenguaje para procesar el
        contenido visual de un documento (tablas, layouts, imágenes) sin requerir
        extracción de texto previa.
\end{description}

% ============================================================
\end{document}